% Editable publishing source. Rebuild without --refresh-tex to preserve edits.
\documentclass[10pt,a4paper]{article}
% Shared editable typography for all five document editions.
\usepackage{fontspec}
\setmainfont{DejaVu Serif}
\setsansfont{DejaVu Sans}
\setmonofont{DejaVu Sans Mono}[Scale=.86]
\usepackage{polyglossia}
\setdefaultlanguage{russian}
\setotherlanguage{english}
\usepackage[a4paper,margin=18mm,top=21mm,bottom=21mm,headheight=15pt]{geometry}
\usepackage{xcolor}
\definecolor{Accent}{HTML}{185A86}
\definecolor{Muted}{HTML}{526573}
\usepackage{graphicx}
\usepackage{longtable,booktabs,array,calc}
\usepackage{fvextra}
\usepackage{pdflscape}
\usepackage{titlesec}
\usepackage{fancyhdr}
\usepackage{enumitem}
\usepackage{etoolbox}
\usepackage{needspace}
\usepackage{tocloft}
\usepackage[unicode,colorlinks=true,linkcolor=Accent,urlcolor=Accent,bookmarksnumbered=false]{hyperref}
\usepackage{bookmark}
\usepackage{microtype}
\AtBeginDocument{\raggedright}
\setcounter{secnumdepth}{0}
\setcounter{tocdepth}{1}
\setlength{\parindent}{0pt}
\setlength{\parskip}{4pt plus 1pt}
\linespread{1.0}
\setlength{\emergencystretch}{3em}
\setlength{\tabcolsep}{4pt}
\renewcommand{\arraystretch}{1.16}
\AtBeginEnvironment{longtable}{\small}
% Measure a compact table or its header and first two data rows. The output
% remains a breakable longtable. Cap the reservation for unusually tall rows.
\newsavebox{\DocumentTableStart}
\newlength{\DocumentTableSpace}
\newcommand{\KeepTableStart}[2]{%
  \setbox\DocumentTableStart=\vbox{\hsize=\linewidth\small #2}%
  \setlength{\DocumentTableSpace}{\dimexpr\ht\DocumentTableStart+\dp\DocumentTableStart
    +\LTpre+\baselineskip+#1\baselineskip*4\relax}%
  \ifdim\DocumentTableSpace>.5\textheight
    \setlength{\DocumentTableSpace}{.5\textheight}%
  \fi
  \Needspace{\DocumentTableSpace}%
}
% Measure short lists once. Small blocks move together; long ones are unboxed
% into the normal page builder and remain breakable, never scaled or clipped.
\newsavebox{\DocumentShortBlock}
\newcommand{\KeepShortBlock}[1]{%
  % A heading's page-space check must not run inside the measuring box: it
  % would inspect the outer page and embed a spurious forced break in the box.
  \par\setbox\DocumentShortBlock=\vbox{\hsize=\linewidth
    \renewcommand{\Needspace}[1]{}#1\par}%
  \ifdim\dimexpr\ht\DocumentShortBlock+\dp\DocumentShortBlock\relax<.32\textheight
    \Needspace{\dimexpr\ht\DocumentShortBlock+\dp\DocumentShortBlock+\medskipamount\relax}%
  \else
    \Needspace{5\baselineskip}%
  \fi
  % Unboxing does not add normal inter-paragraph baseline glue at the leading
  % edge. Give the measured block standard vertical separation from prior text.
  \addvspace{\medskipamount}%
  \unvbox\DocumentShortBlock
}
\preto\section{\Needspace{5\baselineskip}}
\preto\subsection{\Needspace{5\baselineskip}}
\preto\subsubsection{\Needspace{4\baselineskip}}
\clubpenalty=10000
\widowpenalty=10000
\setlist{leftmargin=1.5em,itemsep=2pt,topsep=3pt}
\providecommand{\tightlist}{\setlength{\itemsep}{1pt}\setlength{\parskip}{0pt}}
\titleformat{\section}{\large\sffamily\bfseries\color{Accent}}{}{0pt}{}
\titleformat{\subsection}{\normalsize\sffamily\bfseries\color{Accent}}{}{0pt}{}
\titleformat{\subsubsection}{\normalsize\sffamily\bfseries\color{Accent}}{}{0pt}{}
\titlespacing*{\section}{0pt}{18pt}{8pt}
\titlespacing*{\subsection}{0pt}{13pt}{6pt}
\titlespacing*{\subsubsection}{0pt}{10pt}{5pt}
\pagestyle{fancy}
\fancyhf{}
\fancyfoot[C]{\small\color{Muted}\thepage}
\renewcommand{\headrulewidth}{.35pt}
\renewcommand{\footrulewidth}{0pt}
% A single compact opening page, with profile-controlled TOC depth. Reserve a
% wider page-number column for the multi-page BRD and example editions.
\setlength{\cftbeforesecskip}{5pt}
\setlength{\cftbeforesubsecskip}{2pt}
\renewcommand{\cftsecfont}{\small\sffamily\bfseries}
\renewcommand{\cftsubsecfont}{\small\sffamily}
\renewcommand{\cftsecpagefont}{\small}
\renewcommand{\cftsubsecpagefont}{\small}
\cftsetpnumwidth{2em}
\cftsetrmarg{3em}
\newcommand{\DocumentTitle}[3]{%
  \thispagestyle{plain}
  {\sffamily\bfseries\color{Accent}\fontsize{23}{29}\selectfont #1\par}
  \vspace{3mm}
  {\sffamily\small\color{Muted}#2\par}
  \ifstrempty{#3}{}{\vspace{4mm}#3\par}%
  \begingroup
  \renewcommand{\contentsname}{Содержание}%
  \tableofcontents
  \endgroup
  \vspace{5mm}
  {\color{Accent}\rule{\linewidth}{.8pt}\par}
  \clearpage
}
% Compact standalone profile: keep the title and body on the same opening page.
% It intentionally has no table of contents or forced title-page break.
\newcommand{\CompactDocumentTitle}[3]{%
  \thispagestyle{plain}
  {\sffamily\bfseries\color{Accent}\fontsize{23}{29}\selectfont #1\par}
  \vspace{3mm}
  {\sffamily\small\color{Muted}#2\par}
  \ifstrempty{#3}{}{\vspace{4mm}#3\par}%
  \vspace{5mm}
  {\color{Accent}\rule{\linewidth}{.8pt}\par}
  \vspace{4mm}
}

\hypersetup{pdftitle={Приложение А.2. Эталонный Markdown-шаблон BRD},pdfauthor={},pdfsubject={Шаблон BRD}}
\fancyhead[L]{\small SolutionArchitectAI}
\fancyhead[R]{\small Шаблон BRD}
\setcounter{tocdepth}{1}
\begin{document}
\DocumentTitle{Приложение А.2. Эталонный Markdown-шаблон BRD}{Шаблон BRD}{}
\section{Бизнес-требования (BRD): {[}Название проекта / системы{]}}\label{ux431ux438ux437ux43dux435ux441-ux442ux440ux435ux431ux43eux432ux430ux43dux438ux44f-brd-ux43dux430ux437ux432ux430ux43dux438ux435-ux43fux440ux43eux435ux43aux442ux430--ux441ux438ux441ux442ux435ux43cux44b}

\subsection{0. Метаданные документа}\label{0-ux43cux435ux442ux430ux434ux430ux43dux43dux44bux435-ux434ux43eux43aux443ux43cux435ux43dux442ux430}

Идентификатор документа: BRD-{[}Код проекта{]}

Версия / Редакция: v1.0 (Обязательно менять при любых правках для корректного анализа Diff)

Дата утверждения: ГГГГ-ММ-ДД

Автор (Бизнес-аналитик): {[}ФИО{]}

Заказчик (Владелец продукта): {[}ФИО{]}

Связанный корпоративный контекст: {[}Ссылки на стандарты безопасности, архитектурные политики или действующие исключения{]}

\subsection{1. Цели и назначение системы}\label{1-ux446ux435ux43bux438-ux438-ux43dux430ux437ux43dux430ux447ux435ux43dux438ux435-ux441ux438ux441ux442ux435ux43cux44b}

\subsubsection{1.1. Бизнес-проблема и предпосылки}\label{11-ux431ux438ux437ux43dux435ux441-ux43fux440ux43eux431ux43bux435ux43cux430-ux438-ux43fux440ux435ux434ux43fux43eux441ux44bux43bux43aux438}

{[}Краткое описание решаемой проблемы и причин доработки/создания системы.{]}

\subsubsection{1.2. Бизнес-цели}\label{12-ux431ux438ux437ux43dux435ux441-ux446ux435ux43bux438}

GOAL-01: {[}Измеримая цель, например: сократить время обработки заявки с 2 дней до 15 минут.{]}

GOAL-02: {[}Вторая цель.{]}

\subsubsection{1.3. Границы ответственности (Scope)}\label{13-ux433ux440ux430ux43dux438ux446ux44b-ux43eux442ux432ux435ux442ux441ux442ux432ux435ux43dux43dux43eux441ux442ux438-scope}

Входит в границы проектируемой системы (In-Scope):

{[}Функция 1, например: приём и валидация заявок клиентов.{]}

{[}Функция 2, например: расчёт скорингового балла.{]}

Остаётся за пределами системы (Out-of-Scope / Внешние системы):

{[}Функция 3, например: фактическое списание средств (выполняется в АБС).{]}

{[}Функция 4, например: отправка SMS-уведомлений (выполняется через корпоративный шлюз).{]}

\KeepTableStart{1}{%
\begin{tabular}{@{}
  >{\raggedright\arraybackslash}p{(\columnwidth - 6\tabcolsep) * \real{0.2400}}
  >{\raggedright\arraybackslash}p{(\columnwidth - 6\tabcolsep) * \real{0.2400}}
  >{\raggedright\arraybackslash}p{(\columnwidth - 6\tabcolsep) * \real{0.3000}}
  >{\raggedright\arraybackslash}p{(\columnwidth - 6\tabcolsep) * \real{0.2200}}@{}}
\toprule\noalign{}
\begin{minipage}[b]{\linewidth}\raggedright
ID роли
\end{minipage} & \begin{minipage}[b]{\linewidth}\raggedright
Название роли
\end{minipage} & \begin{minipage}[b]{\linewidth}\raggedright
Тип участника (Человек / Внешняя система)
\end{minipage} & \begin{minipage}[b]{\linewidth}\raggedright
Описание задач и прав в системе
\end{minipage} \\
\midrule\noalign{}
ROLE-01 & Клиент (Физлицо) & Внешний пользователь & Подаёт заявку через мобильное приложение, просматривает статус. \\
ROLE-02 & Менеджер верификации & Внутренний сотрудник & Проверяет заявки, попавшие на ручной разбор, утверждает или отклоняет их. \\
ROLE-03 & Корпоративная АБС & Внешняя система & Хранит счета клиентов, принимает команды на открытие договора. \\

\end{tabular}%
}
\subsection{2. Участники и бизнес-роли}\label{2-ux443ux447ux430ux441ux442ux43dux438ux43aux438-ux438-ux431ux438ux437ux43dux435ux441-ux440ux43eux43bux438}

\begin{longtable}[]{@{}
  >{\raggedright\arraybackslash}p{(\columnwidth - 6\tabcolsep) * \real{0.2400}}
  >{\raggedright\arraybackslash}p{(\columnwidth - 6\tabcolsep) * \real{0.2400}}
  >{\raggedright\arraybackslash}p{(\columnwidth - 6\tabcolsep) * \real{0.3000}}
  >{\raggedright\arraybackslash}p{(\columnwidth - 6\tabcolsep) * \real{0.2200}}@{}}
\toprule\noalign{}
\begin{minipage}[b]{\linewidth}\raggedright
ID роли
\end{minipage} & \begin{minipage}[b]{\linewidth}\raggedright
Название роли
\end{minipage} & \begin{minipage}[b]{\linewidth}\raggedright
Тип участника (Человек / Внешняя система)
\end{minipage} & \begin{minipage}[b]{\linewidth}\raggedright
Описание задач и прав в системе
\end{minipage} \\
\midrule\noalign{}
\endhead
\bottomrule\noalign{}
\endlastfoot
ROLE-01 & Клиент (Физлицо) & Внешний пользователь & Подаёт заявку через мобильное приложение, просматривает статус. \\
ROLE-02 & Менеджер верификации & Внутренний сотрудник & Проверяет заявки, попавшие на ручной разбор, утверждает или отклоняет их. \\
ROLE-03 & Корпоративная АБС & Внешняя система & Хранит счета клиентов, принимает команды на открытие договора. \\
\end{longtable}

\begin{landscape}

\subsection{3. Функциональные требования (FR)}\label{3-ux444ux443ux43dux43aux446ux438ux43eux43dux430ux43bux44cux43dux44bux435-ux442ux440ux435ux431ux43eux432ux430ux43dux438ux44f-fr}

\subsubsection{3.1. {[}Название бизнес-процесса, например: Подача заявки{]}}\label{31-ux43dux430ux437ux432ux430ux43dux438ux435-ux431ux438ux437ux43dux435ux441-ux43fux440ux43eux446ux435ux441ux441ux430-ux43dux430ux43fux440ux438ux43cux435ux440-ux43fux43eux434ux430ux447ux430-ux437ux430ux44fux432ux43aux438}

FREQ-01: Система должна принимать от ROLE-01 заполненную анкету на кредит.

FREQ-02: При получении анкеты система должна запрашивать кредитную историю клиента из внешней системы {[}Название системы{]}.

FREQ-03: Если сумма кредита превышает 10 000 руб., система должна направлять заявку на ручное согласование ROLE-02.

\KeepTableStart{1}{%
\begin{tabular}{@{}
  >{\raggedright\arraybackslash}p{(\columnwidth - 12\tabcolsep) * \real{0.0800}}
  >{\raggedright\arraybackslash}p{(\columnwidth - 12\tabcolsep) * \real{0.1400}}
  >{\raggedright\arraybackslash}p{(\columnwidth - 12\tabcolsep) * \real{0.1500}}
  >{\raggedright\arraybackslash}p{(\columnwidth - 12\tabcolsep) * \real{0.1400}}
  >{\raggedright\arraybackslash}p{(\columnwidth - 12\tabcolsep) * \real{0.1800}}
  >{\raggedright\arraybackslash}p{(\columnwidth - 12\tabcolsep) * \real{0.1600}}
  >{\raggedright\arraybackslash}p{(\columnwidth - 12\tabcolsep) * \real{0.1500}}@{}}
\toprule\noalign{}
\begin{minipage}[b]{\linewidth}\raggedright
ID данных
\end{minipage} & \begin{minipage}[b]{\linewidth}\raggedright
Наименование набора данных
\end{minipage} & \begin{minipage}[b]{\linewidth}\raggedright
Источник (Отправитель)
\end{minipage} & \begin{minipage}[b]{\linewidth}\raggedright
Получатель
\end{minipage} & \begin{minipage}[b]{\linewidth}\raggedright
Состав ключевых атрибутов
\end{minipage} & \begin{minipage}[b]{\linewidth}\raggedright
Класс конфиденциальности
\end{minipage} & \begin{minipage}[b]{\linewidth}\raggedright
Периодичность и условия передачи
\end{minipage} \\
\midrule\noalign{}
DATA-01 & Анкета клиента & ROLE-01 (Клиент) & Проектируемая система & ФИО, паспортные данные, сумма, срок & Персональные данные (ПДн) & По событию (в момент отправки заявки) \\
DATA-02 & Кредитный отчёт & Бюро кредитных историй (БКИ) & Проектируемая система & Скоринг-балл, список открытых кредитов & Конфиденциально & По запросу (синхронно при проверке заявки) \\
DATA-03 & Реестр выданных кредитов & Проектируемая система & ROLE-03 (АБС) & ID заявки, ID клиента, сумма, ставка & Коммерческая тайна & Пакетно, 1 раз в сутки в 23:00 \\

\end{tabular}%
}
\subsection{4. Описание данных и потоков (Data Requirements)}\label{4-ux43eux43fux438ux441ux430ux43dux438ux435-ux434ux430ux43dux43dux44bux445-ux438-ux43fux43eux442ux43eux43aux43eux432-data-requirements}

\begin{longtable}[]{@{}
  >{\raggedright\arraybackslash}p{(\columnwidth - 12\tabcolsep) * \real{0.0800}}
  >{\raggedright\arraybackslash}p{(\columnwidth - 12\tabcolsep) * \real{0.1400}}
  >{\raggedright\arraybackslash}p{(\columnwidth - 12\tabcolsep) * \real{0.1500}}
  >{\raggedright\arraybackslash}p{(\columnwidth - 12\tabcolsep) * \real{0.1400}}
  >{\raggedright\arraybackslash}p{(\columnwidth - 12\tabcolsep) * \real{0.1800}}
  >{\raggedright\arraybackslash}p{(\columnwidth - 12\tabcolsep) * \real{0.1600}}
  >{\raggedright\arraybackslash}p{(\columnwidth - 12\tabcolsep) * \real{0.1500}}@{}}
\toprule\noalign{}
\begin{minipage}[b]{\linewidth}\raggedright
ID данных
\end{minipage} & \begin{minipage}[b]{\linewidth}\raggedright
Наименование набора данных
\end{minipage} & \begin{minipage}[b]{\linewidth}\raggedright
Источник (Отправитель)
\end{minipage} & \begin{minipage}[b]{\linewidth}\raggedright
Получатель
\end{minipage} & \begin{minipage}[b]{\linewidth}\raggedright
Состав ключевых атрибутов
\end{minipage} & \begin{minipage}[b]{\linewidth}\raggedright
Класс конфиденциальности
\end{minipage} & \begin{minipage}[b]{\linewidth}\raggedright
Периодичность и условия передачи
\end{minipage} \\
\midrule\noalign{}
\endhead
\bottomrule\noalign{}
\endlastfoot
DATA-01 & Анкета клиента & ROLE-01 (Клиент) & Проектируемая система & ФИО, паспортные данные, сумма, срок & Персональные данные (ПДн) & По событию (в момент отправки заявки) \\
DATA-02 & Кредитный отчёт & Бюро кредитных историй (БКИ) & Проектируемая система & Скоринг-балл, список открытых кредитов & Конфиденциально & По запросу (синхронно при проверке заявки) \\
DATA-03 & Реестр выданных кредитов & Проектируемая система & ROLE-03 (АБС) & ID заявки, ID клиента, сумма, ставка & Коммерческая тайна & Пакетно, 1 раз в сутки в 23:00 \\
\end{longtable}

\end{landscape}

\subsection{5. Нефункциональные требования (NFR)}\label{5-ux43dux435ux444ux443ux43dux43aux446ux438ux43eux43dux430ux43bux44cux43dux44bux435-ux442ux440ux435ux431ux43eux432ux430ux43dux438ux44f-nfr}

NFR-01 (Доступность / SLA): Система должна быть доступна в режиме 24/7 с коэффициентом доступности не ниже 99.99\%.

NFR-02 (Производительность и нагрузка): Максимальное время ответа на синхронный запрос FREQ-01 не должно превышать 2 секунд при пиковой нагрузке 100 RPS.

NFR-03 (Окно обслуживания и фоновые задачи): Пакетная выгрузка DATA-03 должна выполняться строго в ночное время, чтобы не создавать нагрузку на пользовательские операции.

NFR-04 (Безопасность и защита каналов): Передача DATA-01 и DATA-02 по внешним сетям должна осуществляться исключительно по шифрованным каналам с взаимной аутентификацией (mTLS / ГОСТ).

NFR-05 (Хранение и аудит): Все действия ROLE-02 по изменению статуса заявки должны логироваться и храниться не менее 5 лет.

\subsection{6. Архитектурно значимые ограничения и допущения}\label{6-ux430ux440ux445ux438ux442ux435ux43aux442ux443ux440ux43dux43e-ux437ux43dux430ux447ux438ux43cux44bux435-ux43eux433ux440ux430ux43dux438ux447ux435ux43dux438ux44f-ux438-ux434ux43eux43fux443ux449ux435ux43dux438ux44f}

\subsubsection{6.1. Ограничения (Constraints)}\label{61-ux43eux433ux440ux430ux43dux438ux447ux435ux43dux438ux44f-constraints}

CONST-01: Интеграция с ROLE-03 (АБС) допускается только через корпоративную шину данных (ESB / Kafka); прямое подключение к БД АБС запрещено.

CONST-02: Развёртывание компонентов системы допускается только в закрытом контуре ЦОД заказчика (использование публичных облаков запрещено).

\subsubsection{6.2. Известные исключения из стандартов (если применимо)}\label{62-ux438ux437ux432ux435ux441ux442ux43dux44bux435-ux438ux441ux43aux43bux44eux447ux435ux43dux438ux44f-ux438ux437-ux441ux442ux430ux43dux434ux430ux440ux442ux43eux432-ux435ux441ux43bux438-ux43fux440ux438ux43cux435ux43dux438ux43cux43e}

EXC-01: Для взаимодействия с БКИ допускается использование протокола SOAP вместо корпоративного стандарта REST API на основании согласованного архитектурного исключения №{[}Номер{]} от {[}Дата{]}.

\subsubsection{6.3. Открытые вопросы и пробелы (для уточнения)}\label{63-ux43eux442ux43aux440ux44bux442ux44bux435-ux432ux43eux43fux440ux43eux441ux44b-ux438-ux43fux440ux43eux431ux435ux43bux44b-ux434ux43bux44f-ux443ux442ux43eux447ux43dux435ux43dux438ux44f}

Q-01: Не определён формат файла для пакетной выгрузки DATA-03 (XML или JSON) --- требует уточнения со стороны команды АБС.

\end{document}
